\documentclass[10pt,a4paper]{amsart}
\usepackage[margin=19mm]{geometry}
\usepackage{amsmath,amssymb,mathtools}
\usepackage[scr=rsfs,cal=euler]{mathalfa}
\usepackage{xfrac}
\usepackage[unicode,hidelinks,bookmarks=false]{hyperref}
\newcommand{\ie}{i.e.\ }
\newcommand{\cf}{cf.\ }
\newcommand{\resp}{respectively\ }
\newcommand{\Subsection}{\subsection}
\providecommand{\nobreakdash}{\nobreak\hbox{-}}
\setlength{\emergencystretch}{2em}
\multlinegap0pt
\usepackage{amssymb}
\newtheorem{proposition}{Proposition}
\newtheorem{theorem}{Theorem}
\title{On cyclic invariants of the free associative algebra}
\author{Silvia Boumova, Vesselin Drensky}
\date{Source: arXiv:2602.16202v1 (CC BY 4.0)}
\begin{document}
\maketitle

\noindent\emph{Source and licence.} On cyclic invariants of the free associative algebra. Version arXiv:2602.16202v1 (CC BY 4.0), distributed by arXiv under the Creative Commons Attribution 4.0 International licence, http://creativecommons.org/licenses/by/4.0/, as recorded in the arXiv licence metadata for this exact version and shown on the abstract page https://arxiv.org/abs/2602.16202v1. Full licence notice: \texttt{licenses/arxiv-2602.16202-CC-BY-4.0.txt}.\par\smallskip

\noindent\textit{Editorial scope.} Counted unit: the article's theorem S-invariants (source0.tex lines 565--573). The article's complete original proof of it is delivered with it (source0.tex lines 575--593, one \texttt{proof} environment of 19 lines). No mathematics is written, repaired or re-derived: the statement and the proof are the article's own lines, in the article's own notation, and no retained line is substituted or re-flowed.  Retained uncounted as the source-local context the proof needs: lines 278--285; lines 490--497.  Nothing was omitted from any retained span, so the package carries no omission record.  No retained line contains a citation, so the wrapper carries no bibliography and no bibliography entry.  No retained line contains a figure environment, an image-inclusion command or drawing code, so no image is delivered and the package declares no asset dependency.  Editorial formatting only, in addition: an AMS \texttt{amsart} preamble that restates the article's own declarations and macros used by the retained lines, namely the environments \texttt{proposition}, \texttt{theorem} on separate counters, together with the standard packages those lines need.  The retained environments print in this document as Proposition 1, Theorem 1 in order of appearance, whereas the article prints the same items with the numbering of its own full document.  The complete native source of the article as distributed in the arXiv e-print for this exact version is bundled unchanged as \texttt{sources/194856/source0.tex}.
\begin{proposition}\label{commutative case}
The algebra of invariants ${\mathbb C}[Y_d]^{C_d}$ is generated by $y_0$ and the monomials
$u=y_1^{n_1}\cdots y_{d-1}^{n_{d-1}}$ such that
\[
1\leq n_1+\cdots+n_{d-1}\leq d,n_1+2n_2+\cdots+(d-1)n_{d-1}\equiv 0 \text{ \rm (mod }d)
\]
and which cannot be presented as products of two or more polynomials of this kind.
\end{proposition}

\begin{theorem}\label{description of algebra of invariants}
{\rm (i)} The algebra ${\mathbb C}\langle Y_d\rangle^{C_d}$ has a basis which consists of all monomials
$y_{i_1}\cdots y_{i_n}$ such that $i_1+\cdots+i_n\equiv 0\text{ \rm (mod }d)$.

{\rm (ii)} The set of all monomials from {\rm (i)}
such that for all $m=1,2,\ldots,n-1$ the beginning monomial $y_{i_1}\cdots y_{i_m}$ satisfies $i_1+\cdots+i_m\not\equiv 0\text{ \rm (mod }d)$
forms a set of free generators of ${\mathbb C}\langle Y_d\rangle^{C_d}$.
\end{theorem}

\begin{theorem}\label{cyclic S-invariants}
Let
\[
U=\{u_0=y_0,u_1,\ldots,u_k\},\;u_i=y_1^{n_{i,1}}\cdots y_{d-1}^{n_{i,d-1}},\;i=1,\ldots,k,
\]
be the set of generators of the algebra ${\mathbb C}[Y_d]^{C_d}$ from Proposition \ref{commutative case}.
Then the same monomials considered as elements of ${\mathbb C}\langle Y_d\rangle$ generate the algebra
$({\mathbb C}\langle Y_d\rangle^{C_d},\circ)$.
\end{theorem}

\begin{proof}
By Theorem \ref{description of algebra of invariants} (i) the algebra ${\mathbb C}\langle Y_d\rangle^{C_d}$ has a basis of all monomials
$y_{i_1}\cdots y_{i_n}$ such that $i_1+\cdots+i_n\equiv 0\text{ \rm (mod }d)$.
Hence it is sufficient to show that such monomials belong to the $S$-subalgebra of ${\mathbb C}\langle Y_d\rangle$ generated by the set $U$.
There is a $\sigma\in \text{\rm Sym}_n$ such that
\[
y_{i_1}\cdots y_{i_n}\circ\sigma=y_0^{n_0}y_1^{n_1}\cdots y_{d-1}^{n_{d-1}}
\]
and $n_1+2n_2+\cdots+(d-1)n_{d-1}\equiv 0\text{ \rm (mod }d)$. Working in ${\mathbb C}[Y_d]^{C_d}$,
\[
y_0^{n_0}y_1^{n_1}\cdots y_{d-1}^{n_{d-1}}=u_0^{n_0}u_1^{p_1}\cdots u_k^{p_k}
\]
for suitable $p_1,\ldots,p_k$. Going back to ${\mathbb C}\langle Y_d\rangle^{C_d}$ there is a $\tau\in \text{\rm Sym}_n$ such that
\[
y_0^{n_0}y_1^{n_1}\cdots y_{d-1}^{n_{d-1}}\circ\tau=u_0^{n_0}u_1^{p_1}\cdots u_k^{p_k}.
\]
Hence $y_{i_1}\cdots y_{i_n}\circ\sigma\tau=u_0^{n_0}u_1^{p_1}\cdots u_k^{p_k}$ and this shows that all monomials in the basis of
${\mathbb C}\langle Y_d\rangle^{C_d}$ belong to the $S$-subalgebra of ${\mathbb C}\langle Y_d\rangle^{C_d}$ generated by the set $U$.
\end{proof}

\end{document}
